\chapter{Render trees and reconciliation}
\label{ch:reconcile}

\begin{aside}
The element tree is rebuilt every frame. Naively that means
every cell repaints. The reconciler keeps the rebuild logical
but the paint physical: same content, same cells, no churn.
\end{aside}

\section{The naive approach}

Build a new tree. Throw away the old. Paint from scratch.
Correct but wasteful: every frame is a full repaint.

\section{Reconciliation}

The reconciler compares old and new trees and computes the
minimum set of operations to transform one into the other.
Operations:

\begin{itemize}[leftmargin=*]
  \item Same element kind and props: reuse.
  \item Same kind, changed props: update props.
  \item Different kind: replace.
  \item Subtree added: mount.
  \item Subtree removed: unmount.
\end{itemize}

\maya{} does not do full diff-style reconciliation. Instead,
it rebuilds and paints into the back buffer; the canvas diff
(Chapter~\ref{ch:diffloop}) does the equivalent at cell level.

\section{Why no element-level diff}

\maya{}'s elements are cheap to construct (no allocations
when arena-backed). Painting a redundant element costs the
same as marking it ``unchanged''. The cell-level diff is
where the saving lies.

React reconciles because the DOM is expensive to touch.
\maya{}'s ``DOM'' is a 2-D byte array; touching it costs
one store.

\section{Keys}

The exception: lists with stable identity. A list of items
that reorder need keys so the renderer matches them across
frames:

\begin{lstlisting}
list_for_each(items, [](Item& it) {
    return row(...).key(it.id);
});
\end{lstlisting}

The key tells the runtime ``this is the same row as last
frame'', preserving cursor focus, animation state, and any
per-element scratch.

Without keys, a reorder looks like ``everything changed''
to the focus and animation systems.

\section{Stateful widgets}

A text input has internal state (cursor position, selection,
undo stack). When the tree rebuilds, the widget's state must
survive.

\maya{} keys widget instances by tree path: ``root > column
> input[2]''. Same path on the next frame returns the same
state instance.

\section{Path collisions}

If two children have the same kind and no key, their paths
collide. Re-orders confuse the state mapping. \maya{} warns
in debug builds and asks for explicit keys.

\section{Mounting and unmounting}

A widget that appears (mounted) gets a chance to initialise
(register effects, start animations). A widget that
disappears (unmounted) gets a chance to clean up.

\maya{}'s lifecycle hooks:

\begin{lstlisting}
ctx.on_mount(my_widget, [&] { ... });
ctx.on_unmount(my_widget, [&] { ... });
\end{lstlisting}

The runtime walks the tree, compares with last frame's set,
and calls hooks for new and missing widgets.

\section{Effect lifetimes}

Effects registered inside a widget live as long as the
widget. Unmount drops the effect. This prevents leaks (no
dangling effect firing on a destroyed widget).

\section{Recap}

\begin{recap}
\begin{itemize}[leftmargin=*]
  \item \maya{} rebuilds the element tree per frame.
  \item Cell-level diff (not element-level) saves work.
  \item Keys identify list rows across frames.
  \item Widget state survives via tree-path identity.
  \item Mount/unmount hooks bound effect lifetimes.
\end{itemize}
\end{recap}

\section*{Workshop}

\begin{enumerate}[leftmargin=*]
  \item Build a list that reorders. Compare with and
    without keys; observe what happens to focus and
    animations.
  \item Add a mount/unmount log; watch widgets appear and
    disappear as panels toggle.
  \item Find a state leak by mounting a widget repeatedly
    and watching memory. Fix with proper unmount.
\end{enumerate}
